\documentclass[9pt,a4paper]{article}

\usepackage[a4paper,top=12mm,bottom=12mm,left=15mm,right=15mm]{geometry}
\usepackage{fontspec}
\usepackage{xcolor}
\usepackage{tabularx}
\usepackage{enumitem}
\usepackage[most]{tcolorbox}
\usepackage{microtype}
\usepackage{array}

\IfFontExistsTF{Arial}{
  \setmainfont{Arial}
  \setsansfont{Arial}
}{
  \setmainfont{TeX Gyre Heros}
  \setsansfont{TeX Gyre Heros}
}

\definecolor{AIEPRed}{HTML}{D62027}
\definecolor{AIEPNavy}{HTML}{082B5C}
\definecolor{AIEPInk}{HTML}{182B3A}
\definecolor{AIEPMuted}{HTML}{5F6B7A}
\definecolor{AIEPLine}{HTML}{D8DEE6}
\definecolor{AIEPSoft}{HTML}{F5F2EC}

\pagestyle{empty}
\setlength{\parindent}{0pt}
\setlength{\parskip}{3pt}
\setlist[itemize]{leftmargin=*,topsep=1pt,itemsep=1pt,parsep=0pt}
\setlist[enumerate]{leftmargin=*,topsep=1pt,itemsep=1pt,parsep=0pt}
\renewcommand{\arraystretch}{1.08}

\newtcolorbox{softbox}[1]{
  colback=AIEPSoft,
  colframe=AIEPLine,
  boxrule=0.5pt,
  arc=1mm,
  left=2.5mm,right=2.5mm,top=1.7mm,bottom=1.7mm,
  before skip=3pt,after skip=3pt,
  title={\scriptsize\bfseries\MakeUppercase{#1}},
  coltitle=AIEPNavy,
  attach title to upper={\par\vspace{0.5mm}}
}

\newtcolorbox{accentbox}[1]{
  colback=white,
  colframe=AIEPRed,
  boxrule=0.65pt,
  arc=1mm,
  left=2.5mm,right=2.5mm,top=1.7mm,bottom=1.7mm,
  before skip=3pt,after skip=3pt,
  title={\scriptsize\bfseries\MakeUppercase{#1}},
  coltitle=AIEPRed,
  attach title to upper={\par\vspace{0.5mm}}
}

\begin{document}

{\small\bfseries\textcolor{AIEPRed}{GEOGREEN ESCOLAR · AIEP OSORNO}}\par
\vspace{2mm}
{\Huge\bfseries\color{AIEPNavy}Léeme — Mentoría 4\par}
\vspace{1mm}
{\Large\bfseries\color{AIEPInk}Comunicar y ensayar\par}
{\normalsize\color{AIEPMuted}Material base para el equipo docente de Desarrollo Social\par}

\vspace{4mm}
Esta carpeta reúne los materiales para ejecutar la Mentoría 4 de GeoGreen Escolar. Puede utilizarlos directamente o ajustar su mediación al grupo. La organización, los productos de entrada y la salida obligatoria corresponden a la lógica común del programa y deben preservarse.

\begin{softbox}{Qué hace esta mentoría}
La sesión transforma el trabajo desarrollado por cada equipo en una presentación comprensible, coherente y compartida. Para ello, el grupo:
\begin{enumerate}
  \item recupera la relación entre problema, contexto, solución y evidencia;
  \item construye un relato breve con una idea principal por tramo;
  \item selecciona un apoyo visual que permita explicar o demostrar;
  \item distribuye contenido, acciones y transiciones entre sus seis integrantes;
  \item ensaya, recibe retroalimentación y comprueba una corrección concreta.
\end{enumerate}
\end{softbox}

\begin{accentbox}{Condiciones de ejecución}
\begin{tabular}{@{}>{\bfseries\color{AIEPNavy}\raggedright\arraybackslash}p{39mm}@{\hspace{5mm}}>{\raggedright\arraybackslash}p{\dimexpr\linewidth-44mm\relax}@{}}
\strut Duración & \strut 60 minutos por bloque. \\[1mm]
\strut Organización & \strut 5 equipos de 6 estudiantes por bloque. \\[1mm]
\strut Entradas obligatorias & \strut Problema y contexto M1, solución y recursos M2, y evidencia de avance M3. \\[1mm]
\strut Salida obligatoria & \strut Guion del pitch, apoyo visual, seis roles con tiempos y una corrección comprobada. \\[1mm]
\strut Rol de la mentoría & \strut Estructurar, observar y retroalimentar; la preparación continúa con trabajo del equipo antes de la presentación final. \\
\end{tabular}
\end{accentbox}

\begin{softbox}{Preparación breve}
Revise la presentación e imprima una ficha por equipo —cinco por bloque—. Solicite que cada grupo lleve sus productos de las mentorías anteriores y el recurso visual que desea mostrar. Disponga un cronómetro visible para los ensayos.
\end{softbox}

{\bfseries\color{AIEPNavy}Qué hay en esta carpeta}\par
\begin{itemize}
  \item \textbf{LEEME.pdf} — esta guía de orientación.
  \item \textbf{Mentoría 4 - mapa de la clase.png} — síntesis visual de la sesión, útil para compartir y proyectar.
  \item \textbf{GeoGreen Escolar - resumen del proyecto.png} y \textbf{GeoGreen Escolar - lógica del proyecto.png} — panorama general y secuencia de talleres, mentorías y productos.
  \item \textbf{Documentos/} — planificación docente completa y ficha imprimible para el trabajo de cada equipo.
  \item \textbf{Presentacion/} — apoyo visual para conducir la sesión en PowerPoint.
\end{itemize}

\begin{tcolorbox}[colback=AIEPNavy,colframe=AIEPNavy,coltext=white,boxrule=0pt,arc=1mm,left=3mm,right=3mm,top=1.8mm,bottom=1.8mm,before skip=3pt]
\textbf{Resultado esperado:} cada equipo dispone de una presentación ensayada en la que seis voces construyen una misma historia, la evidencia cumple una función clara y al menos una mejora queda incorporada.
\end{tcolorbox}

\vfill
{\footnotesize\color{AIEPMuted}GeoGreen Escolar · AIEP Osorno — Programación, Electricidad y Electrónica, y Trabajo Social.}

\end{document}
